\begin{exercise}{81}
Let $\mathcal C_0\subseteq\mathcal C$ be a full subcategory. Show that the inclusion
$\mathcal C_0\to\mathcal C$ is an isofibration if $\mathcal C_0$ is closed under
equivalences in $\mathcal C$, i.e., that if $x\in\mathcal C_0$ and $y\in\mathcal C$
is equivalent to $x$, then $y$ is also in $\mathcal C_0$.

\begin{proof}
By Definition~1.2.76, a simplex of $\mathcal C$ belongs to
$\mathcal C_0$ exactly when all its vertices do;
Lemma~1.2.77 ensures that $\mathcal C_0$ is an $\infty$-category.
Every inner horn contains all vertices of its ambient simplex.
Thus, in a lifting problem for $\mathcal C_0\hookrightarrow
\mathcal C$ against an inner horn inclusion, the prescribed
simplex in $\mathcal C$ already lies in $\mathcal C_0$.
This proves that the inclusion is an inner fibration.

Now let $e\colon x\to y$ be an equivalence in $\mathcal C$
with $x\in\mathcal C_0$.
The hypothesis gives $y\in\mathcal C_0$, so fullness puts
$e$ in $\mathcal C_0$.
The two inverse triangles for $e$ from Lemma~1.2.23 have
only $x$ and $y$ as vertices, so they also lie in
$\mathcal C_0$.
The same lemma therefore makes $e$ an equivalence in
$\mathcal C_0$.
This solves every equivalence-lifting problem in
Definition~2.1.4, proving that the inclusion is an isofibration.
\end{proof}
\end{exercise}
